\section{Related Work}
\label{sec:related_work}

\paragraph{Vision-Language Models.}
Inspired by the success of language foundation models such as GPT-3~\cite{brown2020language} and T5~\cite{raffel2020exploring}, a series of work pre-train vision-language models on large-scale image-text datasets~\cite{lei2021less,radford2021learning,jia2021scaling}. Among them, Contrastive Language-Vision Pre-training~\cite{radford2021learning} achieves remarkable performance on various downstream tasks. It concentrates on aligning images and texts to acquire a joint embedding space.
The CLIP model contains an image encoder~\cite{he2016deep,dosovitskiy2020image} and a text encoder~\cite{devlin2018bert}. During pre-training, contrastive learning is performed in which a paired image-text is a positive pair while image and text from different image-text pairs form a negative pair. 
% For inference, each class is transformed into a sentence using a template such as ``A photo of \{classname\}.'' Then 
For inference, the closest text embedding for the image is chosen as the prediction. Vision-language models can give zero-shot predictions on unseen tasks with a robust zero-shot transfer ability on various downstream tasks.

\paragraph{Continual Learning Methods.} Most existing continual learning methods can be categorized into four groups: parameter expansion, memory replay, distillation loss, and parameter regularization. Parameter expansion methods such as DyTox~\cite{douillard2022dytox} and DEN~\cite{yoon2017lifelong} 
% and so on~\cite{serra2018overcoming,li2019learn,hung2019compacting,golkar2019continual}
introduce new parameters for new tasks. As we want to achieve a more powerful CLIP model at the end of CL, we do not change the architecture of the CLIP model. Memory replay methods~\cite{shin2017continual,lopez2017gradient,prabhu2020gdumb,lavda2018continual,ostapenko2022continual} including iCaRL~\cite{rebuffi2017icarl} and SER~\cite{isele2018selective} keep a memory for exemplars from previously learned tasks. However, pre-training datasets are too large for choosing exemplars or may not be available at downstream training, and downstream data are not good exemplars for preserving the pre-training knowledge. Distillation loss such as LwF~\cite{li2017learning}, LwM~\cite{dhar2019learning}, LwF-VR~\cite{ding2022don}, and PodNet~\cite{douillard2020podnet} aligns current output space with previous ones, whereas distillation based on current tasks are not strong enough to maintain foundational knowledge. For the CLIP model, we find that general images, even if never seen by the model, plus sentences with enough semantics, can be a good ``replay'' for the distillation loss. The parameter regularization loss restricts the flexibility of model parameters by training loss~\cite{kirkpatrick2017overcoming,zenke2017continual,aljundi2018memory} or weight averaging~\cite{lee2017overcoming,wortsman2022robust}. Although this type of strategy performs badly compared to other ways in previous research~\cite{van2019three,belouadah2021comprehensive}, we found it helpful for CLIP continual learning. The limited parameter space prevents the model from diverging significantly from its original state.

\paragraph{Vision-Language Models for Downstream Tasks.}
Many works propose different training strategies of vision-language models for better performance on downstream tasks, such as CoOp~\cite{zhou2022learning}, CLIP-Adapter~\cite{gao2021clip} and WiSE-FT~\cite{wortsman2022robust}. However, very few attempts at continual learning exist. While \cite{srinivasan2022climb,wang2022continual} focus on CL in the pretrain of VL and \cite{ni2023continual} trains VL from scratch using a small dataset, our problem setting differs as we address CL in downstream tasks with a pretrained VL. Currently, most VL models are trained directly on an accessible large dataset, yet integrating knowledge continuously into a pretrained VL is a practical necessity. None of these studies address the zero-shot transfer degradation phenomenon, instead focusing on the traditional forgetting of learned knowledge. Recently, Thengane~\etal~\cite{thengane2022clip} shows CLIP zero-shot prediction achieves state-of-the-art performance in CL settings even without any training. LwF-VR~\cite{ding2022don} is a modified LwF method for the CLIP model where random-generated sentences are used for distillation loss. However, it only considers the feature space, and the distillation with random sentences cannot protect the vision backbone. Differently, we re-examine what should be used for distillation in the feature space and combine the parameter space weight ensemble to provide better performance for the vision-language model continual learning.